En los últimos años, se está produciendo una digitalización en la cual, muchas de las tareas que tradicionalmente se han tenido que realizar en persona se están realizando a través de Internet, como trámites bancarios, gestión de documentos personales, compras, mensajería… Toda operación que sea realizada a través de la red, tiene un riesgo asociado a ella: las intrusiones. Una intrusión se define como un acceso no autorizado a una red, dispositivo o sistema con el objetivo del beneficio del atacante. Tanto usuarios particulares como empresas, pueden ser víctimas de estas intrusiones, con consecuencias que pueden variar dependiendo de la intencionalidad del atacante pero que en ambos casos pueden ser muy graves.

A nivel particular, las consecuencias son muy variadas: captación de datos personales, robo de credenciales, suplantación de identidad o fraude económico. Con que el atacante obtenga acceso a una cuenta, puede ser suficiente para desencadenar una oleada de intrusiones si, por ejemplo, el usuario reutiliza credenciales en distintas plataformas. [FUENTE 1] El Instituto Nacional de Ciberseguridad (INCIBE) gestionó 97348 incidentes de seguridad en 2024, de los cuales el 67,6% afectaron a la ciudadanía. Entre estos, los incidentes más comunes son el uso de Malware con 42136 casos, incluyendo virus y otros tipos de software maliciosos como Ransomware, que cifra los ficheros del usuario y se pide un pago para descifrarlos. Detrás de estos, les sigue el fraude online con más de 38000 incidentes. De estos últimos, el phishing fue el ataque más común, en el que el atacante finge ser entidades legítimas para apropiarse de datos bancarios o personales.

A nivel de empresa, el impacto puede ser incluso mayor si se tienen en cuenta la escala de los activos afectados (que varían desde monetarios hasta los propios empleados) y la dependencia que tienen las empresas en la infraestructura digital para operar.  [FUENTE 2] Un ejemplo reciente es el ataque sufrido por Logaritme, la principal empresa distribuidora de material sanitario en Cataluña. Los atacantes utilizaron un Ransomware y pidieron un rescate de entre 1 y 3 millones de euros para la recuperación de los datos. Se filtraron datos de unos 300 empleados además de la propia caída del sistema de pedidos.

Con este contexto, resulta evidente que además de las medidas preventivas tradicionales es necesaria la utilización de mecanismos de detección temprana de intrusiones. Es aquí donde aparecen los Sistemas de Detección de Intrusiones (IDS), cuyo objetivo es detectar actividad sospechosa o maliciosa en el tráfico de red. Cuando un IDS detecta comportamiento anómalo, envía una alerta a un equipo de seguridad para la revisión de este tráfico sospechoso. El principal factor limitante de estos es que muchos IDS se basan en reglas de ataques conocidos, por lo que una variación en el ataque o un ataque nuevo puede pasar desapercibido. Esto se puede paliar aplicando un modelo de aprendizaje automático que sea capaz de predecir el ataque a partir de la información de los paquetes del tráfico de red.

En este Trabajo de Fin de Grado, se aplicarán técnicas de aprendizaje automático ya que permite entrenar modelos capaces de aprender patrones a partir de datos y generalizar a situaciones similares. Más concretamente, se abordará la detección de intrusiones como un problema de clasificación, proponiendo un enfoque supervisado basado en Random Forest, entrenado y evaluado sobre un conjunto de datos de referencia.
